The full cochain reconstruction has a precise homotopy consequence.
A space is of finite type here when its integral homology is finitely
generated in each degree; nilpotence means that its fundamental group
is nilpotent and acts nilpotently on all higher homotopy groups.

\begin{corollary}[Horel's and Mandell's recovery theorems applied to the action]
\label{cor:homotopy-recovery}
For connected nilpotent finite-type simplicial sets $X,Y$, the full
natural divided-power test action of Theorem~\ref{thm:reconstruction}
recovers their integral weak homotopy types and determines
\[
 \mathbb R\Map_{\mathsf{cBRing}}(A_Y,A_X)
 \simeq\mathbb R\Map_{\mathsf{Spaces}}(X,Y),
\]
where $A_X^n=\operatorname{Map}(X_n,\Z)$ has its unique pointwise
binomial operations, and $\mathsf{cBRing}$ is the category of
cosimplicial binomial rings with quasi-isomorphisms as weak
equivalences. Object detection also holds for a finite disjoint union
of such spaces. In the derived category of integral $E_\infty$
cochain algebras, object detection and faithfulness still hold, but
fullness fails: for $X=S^2,Y=S^3$, the cochain homotopy set is
$\widehat\Z$ whereas $[S^2,S^3]$ is a singleton.
\end{corollary}
This combines reconstruction with Horel's Theorem~5.4, p.~13, and
Mandell's Main Theorem and Theorems~0.1--0.2, pp.~213--216
\cite{Horel,Mandell}. Section~\ref{sec:homotopy-recovery} verifies
the binomial structure and the derived adjunction explicitly.

The dessin calculations lead to elementary arithmetic examples.
The classical sextic passport $[3,2,1;2,2,1,1;6]$ is Schneps's
three-member orbit \cite[Section~IV, Example~3, pp.~27--28]{Schneps}.
Section~\ref{sec:sextic} gives its complete normalized functions,
monodromy, fields, and collision parameters. The following theorem
specifies the related radical actions and the ambient groups in
which their images are open.

\begin{theorem}[Exact rank-one images and the scalar Kummer bound]
\label{thm:kummer-main}
Choose compatible roots $\theta_N=2^{1/N}>0$ and
$\zeta_N=e^{2\pi i/N}$. Write an element of the absolute Galois
group of $\Q$ as $(a,b)$ on them, so
$\sigma(\zeta_N)=\zeta_N^a$ and
$\sigma(\theta_N)=\zeta_N^b\theta_N$ at every $N$.
The image in $\widehat\Z\rtimes\widehat\Z^\times$, acting by
$x\mapsto ax+b$, is exactly
\[
 (-1)^b=\chi_2(a),\qquad
 \chi_2(a)=
 \begin{cases}1&a\equiv1,7\pmod8,\\-1&a\equiv3,5\pmod8.
 \end{cases}
\]
For the compatible roots $r_N=\zeta_{2N}\theta_N^{-1}$ of $-1/2$,
the translation is $c=(a-1)/2-b$, and its image is exactly
\[
 (-1)^c=\chi_{-2}(a),\qquad
 \chi_{-2}(a)=
 \begin{cases}1&a\equiv1,3\pmod8,\\-1&a\equiv5,7\pmod8.
 \end{cases}
\]
Each image has adelic index two and finite-level order
$N\varphi(N)/d_N$, where $d_N=2$ if $8\mid N$ and $d_N=1$
otherwise.

More generally, let $K$ be a number field and let
$\xi_1,\ldots,\xi_r\in K^\times$ be independent modulo roots of
unity. Choose a finite set $S$ of finite places containing every
nonzero valuation of the $\xi_j$. In
$U_S=\mathcal O_{K,S}^{\times}/\mu_K$, let $M$ be their span and
$M^{\mathrm{sat}}=(M\otimes\Q)\cap U_S$. If a positive integer $e$ annihilates
$M^{\mathrm{sat}}/M$, put $D=e|\mu_K|$.
The translation image over $K(\mu_\infty)$ contains
$D\widehat\Z^r$, and the image over $K$ in the scalar affine group
has index at most
\[
 [K\cap\Q(\mu_\infty):\Q]D^r.
\]
For $r\ge2$ that image has infinite index in the affine group with
full linear part $\operatorname{GL}_r(\widehat\Z)$.
\end{theorem}
The general openness is classical Kummer theory. Schinzel's abelian
radical obstruction, stated in \cite[Theorem~15, p.~7]{PeruccaSgobba},
and Perucca--Sgobba's uniform degree bounds
\cite[Theorems~1 and~13, pp.~1 and~6--8]{PeruccaSgobba}
are the relevant earlier results. Section~\ref{sec:kummer} proves
the explicit images and the displayed saturation bound, and derives
the diagonal action and indices for three Fermat fibres.
